% Carlitz's proof that there are infinitely many irregular primes.

This chapter records the Carlitz route from Kummer congruences to the
infinitude of irregular primes.  The theorem proved in Lean is
\[
  \{p : \N \mid p\text{ is prime and }p\text{ is not regular}\}
  \quad\text{is infinite}.
\]
The proof is by finite-set escape: given a finite set \(S\) of candidate
irregular primes, construct a new prime \(p\notin S\) dividing the numerator
of a suitable divided Bernoulli number \(B_M/M\), then use Kummer congruences
and Kummer's criterion to show that \(p\) is irregular.

\section{The divided Kummer congruence}

The Carlitz argument needs a Kummer congruence for the divided Bernoulli
numbers \(B_m/m\) with no upper bound on the exponents.

\begin{theorem}[Full divided Kummer congruence]
  \label{thm:full-divided-kummer}
  \lean{KummerCriterion.bernoulli_div_sModEq_of_modEq_full}
  \leanok
Let \(p\) be an odd prime.  Let \(m,n>0\) be even natural numbers such that
\[
  m \equiv n \pmod {p-1}
  \qquad\text{and}\qquad
  (p-1)\nmid n.
\]
Then there exists \(z\in \Z_p\) with
\[
  \frac{B_m}{m} - \frac{B_n}{n}
  =
  pz
  \qquad\text{in }\Q_p.
\]
\end{theorem}

\begin{proof}\leanok
The proof first reduces to \(p\ge 5\).  If \(p=3\), then \(p-1=2\) divides
every even exponent \(n\), contradicting the non-boundary hypothesis
\((p-1)\nmid n\).

For \(p\ge 5\), the proof is a strengthened elementary Voronoi argument.  The
first input is the strong Faulhaber congruence
\[
  \sum_{x=0}^{p-1} x^h - pB_h
  \in h p^2\Z_p
\]
for positive even \(h\) with \((p-1)\nmid h\).  In Lean this is
the strong Faulhaber lemma%
\lean{KummerCriterion.sum_range_pow_sub_p_mul_bernoulli_strong}.
It expands the power sum by mathlib's Faulhaber formula and proves that every
non-leading summand contains the required \(hp^2\)-factor.  The delicate
terms are handled with von Staudt-Clausen denominator control and binomial
divisibility.

The second input is the strong Voronoi congruence
\[
  (a^h-1)\frac{B_h}{h}
  -
  a^{h-1}\sum_{x=0}^{p-1} x^{h-1}\left\lfloor\frac{xa}{p}\right\rfloor
  \in p\Z_p
\]
for every \(a\) prime to \(p\), formalized as
the strong Voronoi theorem%
\lean{KummerCriterion.voronoi_congruence_mod_p_strong}.
This substitutes the strong Faulhaber congruence into a higher-modulus
Voronoi sum identity and cancels the nonzero rational factor \(hp\) in
\(\Q_p\).

Next choose a generator \(g\) of \((\Z/p\Z)^\times\), represented by a natural
number \(a\).  Since \(g\) has order \(p-1\), the condition
\((p-1)\nmid h\) implies \(a^h-1\) is a \(p\)-adic unit.  Dividing the strong
Voronoi congruence by this unit proves
\[
  \frac{B_h}{h}\in \Z_p,
\]
recorded as
the Voronoi integrality lemma%
\lean{KummerCriterion.bernoulli_div_self_mem_padicInt_of_not_sub_one_dvd_voronoi}.

Finally apply the strong Voronoi congruence to \(m\) and \(n\).  The factors
\(a^m\) and \(a^n\) agree because \(m\equiv n\pmod{p-1}\), and the floor sums
are congruent modulo \(p\) because their predecessor exponents are congruent;
this is the floor-sum comparison%
\lean{KummerCriterion.voronoi_floor_sum_sModEq_of_pred_modEq}.
The already-proved \(p\)-integrality of \(B_m/m\) and \(B_n/n\) lets the
comparison be carried out inside \(\Z_p\), and the difference is a multiple
of \(p\) in \(\Q_p\).
\end{proof}

\section{From divided Bernoulli numerators to irregularity}

The finite-set construction produces a prime divisor of the numerator of
\(B_M/M\), not directly a Kummer-range numerator \(B_{2k}\).  The next result
is the bridge from that unbounded witness to the usual definition of an
irregular prime.

\begin{theorem}[Carlitz divisor criterion]
  \label{thm:carlitz-divisor-criterion}
  \lean{KummerCriterion.not_isRegularPrime_iff_exists_dvd_bernoulli_div_self_num}
  \leanok
  \uses{thm:full-divided-kummer, thm:kummer-final}
Let \(q\) be an odd prime.  Then \(q\) is not regular if and only if there is
a positive even integer \(m\) such that
\[
  q \mid \operatorname{num}\left(\frac{B_m}{m}\right).
\]
\end{theorem}

\begin{proof}\leanok
If \(q\) is not regular, Kummer's criterion gives a Kummer-range index \(k\)
with
\[
  1\le k,\qquad 2k\le q-3,\qquad q\mid (B_{2k})_{\mathrm{num}}.
\]
Set \(m=2k\).  The range bound implies \(q\nmid m\), so \(m\) is a unit in
\(\Z_q\).  Dividing by this unit preserves numerator divisibility, giving
\[
  q\mid \operatorname{num}\left(\frac{B_m}{m}\right).
\]
The Lean helper is
the divisibility-through-division lemma%
\lean{KummerCriterion.dvd_bernoulli_div_self_num_of_dvd_bernoulli_num}.

Conversely assume \(q\mid\operatorname{num}(B_m/m)\) for some positive even
\(m\).  Von Staudt-Clausen excludes the boundary:
\[
  (q-1)\nmid m,
\]
because if \(q-1\mid m\), then \(q\) cannot divide the reduced numerator of
\(B_m/m\).  This is
the boundary-exclusion lemma%
\lean{KummerCriterion.sub_one_not_dvd_of_dvd_num_bernoulli_div_self}.
Let \(m'\) be the least nonnegative residue of \(m\) modulo \(q-1\).  The
boundary exclusion makes \(m'>0\), and parity gives that \(m'\) is even.  It
also satisfies \(m'<q-1\), hence
\[
  1\le m'/2,\qquad 2(m'/2)\le q-3.
\]

Apply \cref{thm:full-divided-kummer} to \(m\) and \(m'\).  Since
\(B_m/m\equiv B_{m'}/m'\pmod q\), the numerator divisibility transfers from
\(B_m/m\) to \(B_{m'}/m'\), and then to \(B_{m'}\).  The transfer is formalized
by
the \(p\)-adic numerator-transfer lemma%
\lean{KummerCriterion.dvd_bernoulli_num_of_padic_congruent_residue}.
Now \(m'=2(m'/2)\) lies in Kummer's range, so Kummer's criterion gives that
\(q\) is not regular.
\end{proof}

\section{The finite-set construction}

The escape argument starts from a finite set \(S\) and builds an even integer
which is divisible by \(q-1\) for every prime \(q\in S\).

\begin{definition}[Divisor-closed base]
  \label{def:irregular-base}
  \lean{KummerCriterion.irregularBase}
  \leanok
For a finite set \(S\subset\N\), set
\[
  C(S)=2\cdot \bigl(\max(3,\sup S)\bigr)!.
\]
\end{definition}

\begin{proof}\leanok
The factor \(2\) makes \(C(S)\) even, and the factorial makes
\[
  q-1\mid C(S)
\]
for every prime \(q\in S\).  The stronger closure property used internally is
that if a prime \(q\) divides \(C(S)\), then \(q-1\mid C(S)\) as well.
\end{proof}

\begin{theorem}[A new numerator prime outside a finite set]
  \label{thm:numerator-prime-for-carlitz-base}
  \lean{KummerCriterion.exists_numerator_prime_for_carlitz_base}
  \leanok
  \uses{def:irregular-base}
For every finite set \(S\subset\N\), there are natural numbers \(M\) and \(p\)
such that
\[
\begin{gathered}
  p\text{ is prime},\qquad p\ne 2,\qquad M>0,\qquad M\text{ is even},\\
  \forall q\in S,\ q\text{ prime}\Longrightarrow q-1\mid M,\\
  p\mid \operatorname{num}\left(\frac{B_M}{M}\right),
  \qquad p\notin S.
\end{gathered}
\]
\end{theorem}

\begin{proof}\leanok
Start with \(C=C(S)\).  The growth theorem
for divided Bernoulli numbers%
\lean{KummerCriterion.exists_large_even_multiple_abs_bernoulli_div_self_gt_one}
chooses \(t\) such that
\[
  \left|\frac{B_{C2^t}}{C2^t}\right|>1.
\]
Set \(M=C2^t\).  This \(M\) is positive and even, and it is still divisible by
\(q-1\) for every prime \(q\in S\).

Since the rational number \(B_M/M\) has absolute value greater than \(1\), its
reduced numerator has absolute value greater than \(1\), hence has a prime
divisor \(p\), by
the numerator-prime lemma%
\lean{KummerCriterion.exists_prime_dvd_num_of_one_lt_abs}.

It remains to prove that this prime was not already in \(S\).  Suppose
\(p\in S\).  Then \(p-1\mid M\) by construction.  But von Staudt-Clausen gives
the exclusion
\[
  p-1\mid M
  \quad\Longrightarrow\quad
  p\nmid \operatorname{num}\left(\frac{B_M}{M}\right),
\]
namely
the von Staudt numerator-exclusion lemma%
\lean{KummerCriterion.not_dvd_num_bernoulli_div_self_of_sub_one_dvd}.
This contradicts the choice of \(p\).  The same exclusion also proves
\(p\ne 2\), since \(p=2\) would make \(p-1=1\) divide \(M\).
\end{proof}

\section{Escaping every finite set}

\begin{theorem}[Finite-set escape]
  \label{thm:irregular-prime-outside-finite-set}
  \lean{KummerCriterion.exists_not_isRegularPrime_not_mem_carlitz}
  \leanok
  \uses{thm:numerator-prime-for-carlitz-base, thm:carlitz-divisor-criterion}
For every finite set \(S\subset\N\), there is a prime \(p\notin S\) such that
\(p\) is not regular.
\end{theorem}

\begin{proof}\leanok
Apply \cref{thm:numerator-prime-for-carlitz-base} to \(S\), obtaining a prime
\(p\notin S\) and a positive even \(M\) with
\[
  p\mid\operatorname{num}\left(\frac{B_M}{M}\right).
\]
The same theorem gives \(p\ne 2\).  Therefore the Carlitz divisor criterion
\cref{thm:carlitz-divisor-criterion} applies and proves that \(p\) is not
regular.
\end{proof}

\section{Infinitude}

\begin{theorem}[Infinitely many irregular primes]
  \label{thm:infinitely-many-irregular-primes}
  \lean{KummerCriterion.infinite_not_isRegularPrime}
  \leanok
  \uses{thm:irregular-prime-outside-finite-set}
The set
\[
  \{p:\N \mid p\text{ is prime and }p\text{ is not regular}\}
\]
is infinite.
\end{theorem}

\begin{proof}\leanok
Assume the set is finite, and let \(S\) be a finite set containing all of its
elements.  By \cref{thm:irregular-prime-outside-finite-set}, there is a prime
\(p\notin S\) which is not regular.  This contradicts the defining property of
\(S\).  Hence no finite set covers all irregular primes, so the set is
infinite.  The finite-set-to-infinite conversion is the elementary lemma
that no finite cover implies infinitude%
\lean{KummerCriterion.infinite_of_forall_finite_set_not_cover}.
\end{proof}

The checked proof contains no bundled Kummer provider or side-condition
package.  The only source bridge into irregularity is Kummer's criterion from
\cref{ch:kummer-final}, and the unbounded congruence used by Carlitz is the
public theorem \cref{thm:full-divided-kummer}.
